Modern smart manufacturing integrates cyber-physical systems, the industrial Internet of Things, and artificial intelligence to create autonomous production environments where unplanned downtime carries severe economic consequences. Despite considerable progress, predictive maintenance research remains fragmented: anomaly detection, remaining useful life estimation, and maintenance scheduling are typically treated as isolated problems; static models degrade under sensor drift; and combinatorial scheduling complexity grows rapidly with fleet size. This thesis proposes the Quantum-Agentic Smart Manufacturing for Anomaly Detection and Predictive Maintenance (QASAMAP) framework, a three-layer integrated architecture that closes the loop from raw sensor streams to constraint-satisfying maintenance decisions within a single real-time system.

The first layer implements spatio-temporal forecasting and supervised learning on a 100{,}000-reading dataset collected from 50 machines over 69 days. A Temporal Graph Convolutional Network (T-GCN) with a cross-correlation lag-graph construction (50 nodes, 538 directed edges) achieves an average $R^2 = 0.7702$ across five sensor channels, outperforming both a tuned standalone LSTM ($R^2 = 0.6853$, $+12.4\%$) and a hybrid LSTM--GCN baseline ($R^2 = 0.7134$, $+8.0\%$). A comprehensive supervised benchmark across XGBoost, LightGBM, CatBoost, HistGradientBoosting, and TabNet on six prediction targets establishes reference performance for anomaly detection, failure-type classification, downtime-risk and maintenance prediction, machine status, and remaining-useful-life regression.

The second layer is a four-tier real-time multi-agent system that consumes Apache Kafka sensor streams, buffers per-machine windows of width $W=10$, and dispatches them through a nine-agent cognitive pipeline (Monitoring, Correlation, Pattern Learning, Diagnosis, Planning, Action, Self-Reflection, Explainability, Reporting) coordinated by a phased orchestrator. Heterogeneous LLM routing (GLM-5.1, Kimi-K2.6, DeepSeek-v4-Pro) matches each agent to a model whose capabilities suit the task. End-to-end validation on 50 streamed machines processes the full pipeline within the 60-second Kafka cycle: the constraints engine forwards 45 of 50 machines to the scheduling layer at the CRITICAL priority level, and the downstream MILP path returns a one-hot- and capacity-feasible schedule for all 45 machines in 36.9~ms.

The third layer formulates the maintenance scheduling problem as both a Mixed-Integer Linear Program (MILP) and a Quadratic Unconstrained Binary Optimisation (QUBO), and benchmarks four solver paradigms on identical instances: classical exact MILP/CBC, the Quantum-Inspired Evolutionary Algorithm (QIEA), Simulated Quantum Annealing (SQA), and the Quantum Approximate Optimisation Algorithm (QAOA) with an XY ring mixer, executed both in classical state-vector simulation and on IBM Heron~r2 hardware (\texttt{ibm\_fez}, 156 qubits). A fair multi-seed scaling benchmark (5 random seeds per fleet size, identical $T$ across solvers, identical pure-cost metric) shows that classical MILP/CBC returns the proven optimum at every tested $N$ from 3 to 150 with sub-second solve times, while QIEA matches the optimum at $N=3$, retains feasibility up to $N=7$ with a 44--52\% gap, and loses one-hot feasibility for $N \geq 10$. SQA fails the one-hot constraint at every tested $N$ owing to the penalty-dominated QUBO landscape. The QAOA XY-mixer construction recovers the MILP optimum at $N=3$ in both simulation and hardware execution by preserving one-hot states by construction; on hardware, the noise-scaling frontier at $N=5$--$10$ machines is set by circuit depth and gate fidelity, not by qubit count.

A complementary complex-constraint experiment with all four constraints enforced --- one-hot machine assignment, slot-window capacity over the $d=4$-hour duration window, the $\rho$-slot risk horizon, and the must-schedule rule --- on streaming-derived problem instances reveals two distinct empirical regimes. At $N=3$, QAOA-XY on \texttt{ibm\_fez} ($p=2$, $4{,}096$ shots, transpiled depth $801$, $644$ CZ gates) recovers the MILP optimum exactly (cost $=300.00$, $10.96\%$ feasibility rate among shots) by post-selecting the lowest-cost feasible bitstring. At $N=5$ the problem is structurally infeasible for MILP/CBC and infeasible for the feasibility-aware metaheuristics (SA, GA, TS, all reaching $4$ constraint violations); QIEA returns a loosely-feasible schedule at cost $3{,}107.80$ ($+165\%$ over MILP best-effort). The QAOA-XY hardware run, by contrast, post-selects a one-hot- and capacity-feasible schedule at cost $1{,}186.02$ (transpiled depth $1{,}226$, $1{,}935$ CZ, $0.24\%$ feasibility rate), only $1.3\%$ above the MILP best-effort and the only solver in the comparison that returns a usable schedule. This is a constraint-native quantum advantage: the XY-ring mixer's structural preservation of one-hot states, combined with post-selection on hardware-measured shots, recovers a feasible schedule in a regime where every classical exact, classical heuristic, and QUBO-based quantum-inspired method tested either fails feasibility or pays a $2.6\times$ cost premium. The contribution of QASAMAP is therefore both (i) an honest, reproducible benchmark across spatio-temporal forecasting, multi-agent cognitive reasoning, and constraint-aware scheduling on a 60-s real-time Kafka budget, and (ii) a concrete identification of the regime --- constraint-tight streaming-derived problems where MILP feasibility breaks down --- in which constraint-native quantum optimisation provides operational value over the best classical alternatives.
